\color{unidarkgreen}{\subsection{Фиксация состояния АТ}}
\color{black}

\begin{enumerate}[label=\arabic{section}.\arabic{subsection}.\arabic*,  align=left, labelsep=4pt, leftmargin=0pt, itemindent=57pt]

\item
Функциональный блок фиксации состояния (ФСАТ) формирует сигналы отключенного положения своей и смежной стороны АТ с учетом количества выключателей на присоединение. Если действие на выключатель
выведено \textit{«Вывод действия на В1(2,3,4)»}, то подразумевается, что выключатель либо не используется, либо выведен в ремонте и его положение принудительно принимается отключенным.

\item
На рисунке \ref{pic_vo:FSAT} приведен алгоритм работы ФСАТ. 

\medskip
\begin{figure}[H]
\centering
\includegraphics[width=0.8\textwidth,height=0.8\textheight,keepaspectratio]{\fbpath/07.02 РЗАТ/041.1811 ФСАТ/img/FSAT.pdf}
\captionof{figure}{Функциональная схема логики ФСАТ}\label{pic_fsat:FSAT}
\end{figure}

\end{enumerate}